\section*{Chapitre \ref{ch:m3:agriculturals-and-softs} --- Agricoles et softs}

\begin{solution}{exo:m3:agriculturals-and-softs:1}
$1.2/14.8 = 8.1\%$.
\end{solution}

\begin{solution}{exo:m3:agriculturals-and-softs:2}
À 4.00~\$~: 7~\% font 28 cents, arrondis à 30~; \omterm{def:m3:agriculturals-and-softs:limit}{limite élargie} à 45. À 6.00~\$~: 42 cents, arrondis à 40~; \omterm{def:m3:agriculturals-and-softs:limit}{limite élargie} à 60.
\end{solution}

\begin{solution}{exo:m3:agriculturals-and-softs:3}
Juillet livre avant la récolte américaine, donc seul le grain déjà en stock peut être livré, soit l'\omterm{def:m3:agriculturals-and-softs:cropyear}{ancienne récolte}. Décembre
livre après la récolte~: \omterm{def:m3:agriculturals-and-softs:cropyear}{nouvelle récolte}.
\end{solution}

\begin{solution}{exo:m3:agriculturals-and-softs:4}
$300 \times 0.022 + 50 \times 0.11 - 10.00 = 6.60 + 5.50 - 10.00 = \$2.10$ par boisseau.
\end{solution}

\begin{solution}{exo:m3:agriculturals-and-softs:5}
Une journée bloquée (30 cents), puis les 30 cents restants tiennent dans la \omterm{def:m3:agriculturals-and-softs:limit}{limite élargie}, et le marché
se traite à sa fair value le jour 2. La perte est de 60 cents, soit 3\,000~\$ par contrat.
\end{solution}

\begin{solution}{exo:m3:agriculturals-and-softs:6}
Les positions acheteuses encombrées doivent vendre sur la même nouvelle, et il reste peu de nouveaux acheteurs~: les ventes forcées ajoutent au
mouvement fondamental et, sur un marché à limite, peuvent bloquer le cours pendant plusieurs jours.
\end{solution}

\begin{solution}{exo:m3:agriculturals-and-softs:7}
Position nette acheteuse de 24.4~\% des positions ouvertes le 22 mars 2022~; position nette vendeuse de 22.8~\% le 9 juillet 2024~;
position nette vendeuse dans 300 semaines sur 559.
\end{solution}

\begin{solution}{exo:m3:agriculturals-and-softs:8}
La corrélation ne donne pas le sens~: les fonds qui suivent les tendances achètent après une hausse des cours, de sorte que leurs
positions peuvent suivre les cours plutôt que les devancer. Le rapport est par ailleurs publié trois jours après les positions
qu'il décrit. Une stratégie a besoin de preuves que les positions prédisent les rendements ultérieurs, testées hors échantillon sur des données
disponibles à l'époque.
\end{solution}

\begin{solution}{pb:m3:agriculturals-and-softs:1}
\textbf{1.} 30 et 45 cents. \textbf{2.} 50~\$ par cent et par contrat (12.50~\$ par tick d'un quart de cent).
\textbf{3.} 4.20~\$, bloqué. \textbf{4.} 3.75~\$, bloqué à la \omterm{def:m3:agriculturals-and-softs:limit}{limite élargie}. \textbf{5.} Jour 3, à 3.70~\$.
\textbf{6.} $50 \times 0.30 \times 5\,000 = \$75\,000$. \textbf{7.} 112\,500~\$. \textbf{8.} 187\,500~\$
(3\,750~\$ par contrat). \textbf{9.} 200\,000~\$. \textbf{10.} Rien pendant le blocage~: il devient un
ordre au marché qui ne s'exécute que lorsque la séance reprend, vers 3.70~\$. \textbf{11.} Les options peuvent continuer à se négocier et
fixer la fair value, avec une prime qui la reflète~; des contrats à terme synthétiques issus d'options révèlent le
vrai prix du contrat bloqué. \textbf{12.} Peut-être d'autres échéances, ou un contrat étranger, avec un \omterm{def:m3:physical-commodity-markets:basisrisk}{risque de base}.
\textbf{13.} Sa couverture vendeuse gagne de la marge en espèces pendant que sa récolte perd de la valeur~: elle reçoit, elle ne verse pas d'appel de marge.
\textbf{14.} Les limites laissent le temps de réunir la marge et freinent les paniques~; leur coût est l'illiquidité. \textbf{15.} Pour que
la marge nécessaire à deux ou trois jours bloqués soit couverte et que la perte reste tolérable, et non selon la volatilité habituelle.
\textbf{16.} Le degré d'encombrement des positions~: les positions encombrées ajoutent aux ventes forcées. \textbf{17.} Elle plafonne le
mouvement quotidien, et non le mouvement total~: le cours continue de bouger les jours suivants. \textbf{18.} La liquidité vaut le plus
quand on ne peut pas négocier~: la capacité de couvrir ailleurs ou de financer la marge est ce qui survit au blocage.
\textbf{19.} \emph{Résultat à retenir~:} 3\,750~\$ par contrat, 187\,500~\$ pour 50 contrats. \textbf{20.} Une
variation quotidienne maximale qui retarde, sans l'annuler, un changement de cours.
\end{solution}

\begin{solution}{iq:m3:agriculturals-and-softs:1}
Stocks de fin de \omterm{def:m3:agriculturals-and-softs:cropyear}{campagne} rapportés à l'utilisation totale~: le coussin reporté sur la récolte suivante. Quand ce coussin est mince, seul un
prix plus élevé peut rationner l'utilisation, et la demande d'un produit de base est inélastique, si bien que le prix monte fortement quand le
\omterm{def:m3:agriculturals-and-softs:wasde}{ratio stocks/utilisation} baisse~; lorsque le coussin est ample, le prix est fixé par le coût du stockage.
\iqlookfor{demande inélastique et réponse convexe du prix.}
\end{solution}

\begin{solution}{iq:m3:agriculturals-and-softs:2}
Un marché bloqué à sa limite quotidienne, avec des ordres d'un seul côté~: les positions ne peuvent pas être clôturées. Le gérer en dimensionnant
les positions pour plusieurs jours de blocage, en couvrant avec des instruments qui restent négociables (options, autres contrats), et en
conservant des liquidités pour la marge.
\iqlookfor{dimensionnement face au blocage, solutions de rechange, financement.}
\end{solution}

\begin{solution}{iq:m3:agriculturals-and-softs:3}
Retenir la date de publication (vendredi), et non la date du rapport (mardi), comme moment où les données sont connues~; normaliser
les positions (part des positions ouvertes, score z calculé sur le passé seul)~; tester sur des années hors échantillon avec
rendements sans chevauchement~; tenir compte des révisions et des changements de catégorie.
\iqlookfor{alignement sur la date de publication et normalisation rigoureuse.}
\end{solution}

\begin{solution}{iq:m3:agriculturals-and-softs:4}
Du soja en entrée, du tourteau et de l'huile en sortie~; sa marge est le crush. Elle se couvre en achetant des contrats à terme sur soja et en vendant
des contrats à terme sur le tourteau et l'huile au ratio du board, ce qui verrouille la marge sur le volume traité prévu.
\iqlookfor{le spread produit-intrant et le ratio.}
\end{solution}

\begin{solution}{iq:m3:agriculturals-and-softs:5}
Les rendements quotidiens sont tronqués à la limite, si bien qu'une VaR à un jour sous-estime le risque. Retenir des scénarios pluriquotidiens avec
le mouvement restant reporté d'un jour sur l'autre, ou mesurer sur l'horizon nécessaire pour sortir~; intégrer les \omterm{def:m3:agriculturals-and-softs:limit}{limites élargies}.
\iqlookfor{l'horizon de liquidité, et non la distribution tronquée à un jour.}
\end{solution}

\begin{solution}{iq:m3:agriculturals-and-softs:6}
Entrées~: consensus et estimations antérieures, positions et limites, liquidité par contrat. Avant midi~: fixer les
limites de risque, annuler les ordres en attente qui ne devraient pas s'exécuter sur l'annonce, pré-calculer les scénarios. À midi~: analyser la
publication, calculer les surprises, router les ordres avec protection de prix, passer aux options ou à d'autres échéances si un
contrat se bloque~; ensuite, réconcilier les données et produire le rapport.
\iqlookfor{analyse, calcul des surprises, protection et solutions de repli.}
\end{solution}
